% Generated draft: requires mathematical and source review.
\documentclass[11pt]{article}
\usepackage[margin=25mm]{geometry}
\usepackage{fontspec}
\setmainfont{Latin Modern Roman}
\usepackage{amsmath,amssymb,mathtools,unicode-math}
\setmathfont{Latin Modern Math}
\usepackage{enumitem,needspace,xurl,hyperref}
\setlength{\parindent}{0pt}
\setlength{\parskip}{5pt}
\setlength{\emergencystretch}{3em}
\newcommand{\N}{\mathbb{N}}
\newcommand{\Z}{\mathbb{Z}}
\newcommand{\Q}{\mathbb{Q}}
\newcommand{\R}{\mathbb{R}}
\newcommand{\C}{\mathbb{C}}
\newcommand{\F}{\mathbb{F}}
\newcommand{\A}{\mathbb{A}}
\newcommand{\PP}{\mathbb{P}}
\newcommand{\E}{\mathbb{E}}
\newcommand{\eps}{\varepsilon}
\newcommand{\dd}{\,\mathrm{d}}
\begin{document}
\section*{UnsolvedMath ID: 701501}
\textbf{Source problem code:} catalogue entry
\subsection{Definitions and mathematical statement}
Let \(n\geq 2\) be an integer, let \(F_n\) be a free group of rank \(n\), and let \(H\leq F_n\). Interpret \emph{existentially closed in \(F_n\)} in the usual relative sense: every existential first-order formula in the language of groups, with parameters from \(H\), that holds in \(F_n\) also holds in \(H\). Equivalently, every finite system of equations and inequations
\[
 w_i(x_1,\ldots,x_m;\mathbf h)=1\quad(1\leq i\leq r),\qquad
 v_j(x_1,\ldots,x_m;\mathbf h)\neq 1\quad(1\leq j\leq s),
\]
where \(m,r,s\) are nonnegative integers, \(\mathbf h\) is a finite tuple from \(H\), and the \(w_i,v_j\) are group words with these parameters, has a solution in \(H^m\) whenever it has a solution in \(F_n^m\).

A subgroup \(H\) is a \emph{free factor} of \(F_n\) if there exists a subgroup \(K\leq F_n\) such that the homomorphism \(H*K\to F_n\) induced by the inclusions is an isomorphism. The subgroup \(K\) may be trivial.

\textbf{Question (A. G. Myasnikov and V. A. Roman'kov, Kourovka Notebook 18.72).} For every integer \(n\geq 2\) and every subgroup \(H\leq F_n\), does existential closedness of \(H\) in \(F_n\) imply that \(H\) is nonabelian and is a free factor of \(F_n\)? No finite generation assumption is imposed on \(H\).

The supplied record does not define existential closedness explicitly; the formulation above uses full existential formulas with parameters from the subgroup, including inequations, relative to the given ambient group.
\subsection{Short English statement}
Must every subgroup that is existentially closed in a nonabelian free group of finite rank be a nonabelian free factor of that group? Here existential closedness means that every finite system of group equations and inequations with parameters in the subgroup that is solvable in the ambient group is already solvable in the subgroup.
\subsection{Sources}
\begin{itemize}
\item Supplied UnsolvedMath local registry record: permanent ID \(701501\), exact source code \texttt{catalogue entry}, title \emph{Kourovka Notebook Problem 18.72}; upstream revision \texttt{d371ee917ac9647fbb06cd6fc529deec6f3dda82}.
\item Supplied source attribution: \emph{The Kourovka Notebook (Unsolved Problems in Group Theory)}, 18th issue (2014), problem 18.72; also described in the record as the 21st issue edition (2026), \texttt{arXiv:1401.0300v48}. Proposed by A. G. Myasnikov and V. A. Roman'kov in 2014. The collection description names E. I. Khukhro and V. D. Mazurov as editors of Issues 1--21 (1965--2026).
\item The registry labels the problem \emph{open}. This is a supplied status label, not an independently verified literature assessment. The cited source has not been independently consulted for this draft.
\end{itemize}
\end{document}
